\subsection{Mode-pair compensation in the fully weighted theta correlation}
\label{sec:modepaircompensation}

The known fixed-difference sign obstruction does not settle the
fully weighted Fourier correlation. It is therefore natural to
retain the actual arithmetic summands of the theta kernel and
ask whether their pairwise correlation functions possess
nonnegative Fourier transforms. The answer is negative already
at the diagonal level, at \emph{every odd Laguerre order}.
The calculation also isolates an exact modular cancellation
that any constructive replacement must retain.

Let \(q_m=\pi m^2\), and for \(r\ge0\) put
\[
\phi_m(r)=q_m(2q_me^{4r}-3)e^{5r-q_me^{4r}},
\quad
\nu_{m,t}(u)=\tfrac12e^{tu^2}\phi_m(|u|).
\]
For every real \(t\), these functions are strictly positive,
even and integrable; their sum equals the whole-line theta
measure \(\nu_t=\frac12e^{tu^2}\Phi(u)\). Write
\[
h_{m,t}(z)=\int_{\mathbb R}\nu_{m,t}(u)e^{izu}\,du,
\qquad H_t(z)=\sum_{m\ge1}h_{m,t}(z).
\]
For \(n\ge0\) define the ordered correlation
\begin{equation}
G^{mk}_{n,t}(w)=\frac{1}{2(2n)!}
\int_{\mathbb R}b^{2n}
\nu_{m,t}\!\left(\frac{w+b}{2}\right)
\nu_{k,t}\!\left(\frac{w-b}{2}\right)\,db.
\label{eq:modepairdef}
\end{equation}
Every \(G^{mk}_{n,t}\) is pointwise nonnegative, and
\(G_{n,t}=\sum_{m,k\ge1}G^{mk}_{n,t}\) converges in
\(L^1(\mathbb R)\). The Jacobian of
\((u,v)\mapsto(w=u+v,b=u-v)\) gives the exact identity
\begin{equation}
\widehat G^{mk}_{n,t}(x)
=[y^{2n}]\,h_{m,t}(x+iy)h_{k,t}(x-iy),
\qquad
L_{n,t}(x)=\sum_{m,k\ge1}\widehat G^{mk}_{n,t}(x).
\label{eq:modepairFourier}
\end{equation}

\begin{theorem}[Sign obstruction for theta mode-pair Fourier transforms]
\label{thm:modepairsign}
For the genuine square-lattice Jacobi-theta summands,
define
\[
D_m=\phi'_m(0)
=q_me^{-q_m}(-8q_m^2+30q_m-15).
\]
Then
\begin{equation}
D_1>0,\quad D_m<0\ (m\ge2),\quad
\sum_{m\ge1}D_m=0.
\label{eq:modepairDs}
\end{equation}
For fixed \(m,k,n\) and finite real heat time \(t\),
as real \(|x|\to\infty\),
\begin{equation}
\boxed{\;
\widehat G^{mk}_{n,t}(x)
=\frac{(-1)^n(n+1)D_mD_k}{x^{2n+4}}
+O_{m,k,n,t}(|x|^{-2n-6}).
\;}
\label{eq:modepairasymptotic}
\end{equation}
Consequently, for every odd \(n\ge1\), each diagonal
term \(\widehat G^{mm}_{n,t}(x)\) is strictly negative
for all sufficiently large \(|x|\), \emph{despite}
\(G^{mm}_{n,t}(w)\ge0\). At those orders the
\(m=1\) versus \(k\ge2\) mixed terms have a
strictly positive leading Fourier coefficient.

For every finite \(M\ge1\), put
\(H_{M,t}=\sum_{m\le M}h_{m,t}\) and
\(D_M=\sum_{m\le M}D_m>0\).
Then, for each fixed odd \(n\ge1\),
\begin{equation}
\boxed{\;
L_n[H_{M,t}](x)=
-\frac{(n+1)D_M^2}{x^{2n+4}}
+O_{M,n,t}(|x|^{-2n-6})<0
\;}
\label{eq:finiteoddorders}
\end{equation}
at sufficiently large real \(|x|\).
This applies even at \(t=\frac12\), where the
\emph{complete} transform \(H_{1/2}\) is known
to have only real zeros and satisfies the
generalized Laguerre inequalities.
\end{theorem}

\begin{proof}
The theta summands are strictly positive for
\(r\ge0\), since \(q_m>\pi>3\).
The derivative recursion
\(P_0(q)=2q-3\),
\(P_1(q)=-8q^2+30q-15\) gives the displayed
formula for \(D_m\).
For \(q\ge12\), writing \(q=12+v\) with \(v\ge0\)
gives
\[
-P_1(q)=8v^2+162v+807>0.
\]
For the first mode, \(3<\pi<22/7\) and
\(P_1\) is strictly decreasing on this interval;
\(P_1(22/7)=13/49>0\). Thus the first two
sign assertions in \eqref{eq:modepairDs} are exact.
The Jacobi modular equation makes the full
theta kernel even, whence
\(\sum_mD_m=\Phi'(0)=0\). Absolute convergence of
the derivative series justifies this step.
In particular \(D_M=-\sum_{m>M}D_m>0\).

Put \(f_{m,t}(r)=e^{tr^2}\phi_m(r)\) on \(r\ge0\).
For each fixed \(t\), \(f_{m,t}\) and all its
derivatives decay superexponentially as
\(r\to+\infty\); furthermore
\(f'_{m,t}(0)=D_m\).
Repeated integration by parts in
\[
h_{m,t}(z)=\int_0^\infty
f_{m,t}(r)\cos(zr)\,dr
\]
gives, uniformly as \(|\Re z|\to\infty\) in
any fixed horizontal strip,
\begin{equation}
h_{m,t}(z)=-D_mz^{-2}+O_{m,t}(|z|^{-4}).
\label{eq:modepairhalfline}
\end{equation}
Repeated integration by parts to higher
order, followed by Cauchy estimates on a
fixed auxiliary strip, controls any finite
number of derivatives of the remainder.

Multiplying the leading expressions
\(h_{m,t}(x+iy)\) and \(h_{k,t}(x-iy)\)
gives
\[
\frac{D_mD_k}{(x+iy)^2(x-iy)^2}
=\frac{D_mD_k}{(x^2+y^2)^2}.
\]
The coefficient of \(y^{2n}\) is
\((-1)^n(n+1)D_mD_k/x^{2n+4}\).
The next terms and their derivatives
contribute \(O(|x|^{-2n-6})\).
Equation \eqref{eq:modepairFourier} proves
\eqref{eq:modepairasymptotic}. Summing over
a finite rectangle \(m,k\le M\) proves
\eqref{eq:finiteoddorders} and the asserted
signs.
\end{proof}

The complete leading coefficients vanish for
\emph{every} \(n\), not because individual
pairs are Fourier-positive, but because
\[
\sum_{m,k}D_mD_k
=\left(\sum_mD_m\right)^2=0.
\]
More specifically, write
\(D_1=A>0\) and
\(\sum_{m\ge2}D_m=-A\). The diagonal
first-mode and tail--tail blocks together
give \(2A^2\) in this quadratic coefficient;
the two mixed blocks give \(-2A^2\).
For odd \(n\) the former contributions
have negative leading Fourier sign and the
mixed blocks compensate them positively.
This exact equality is the first member of
the infinite family of cancellations involving
all higher odd endpoint jets.

The theorem rules out a termwise Bochner
argument on the mode-pair densities, including
at the heat time \(t=1/2\) where the global
Laguerre inequalities already hold. It does
\emph{not} rule out a nontrivial compensated
infinite regrouping. It also does not prove
the weighted Fourier inequality at \(t=0\):
that all-order inequality is the unresolved
real-zero condition.

A constructive next test should therefore
derive a signed, modularly compensated
\emph{entire-series} representation for
\(\widehat G_{n,t}\), with quantitative control
of the remainder and a positive lower
comparison valid for all \(n,t,x\).
An absolute truncation-error estimate alone
will not suffice, because the Fourier
transform can approach zero at high frequency.
The correlation methods of
Csordas~\cite{Cs,CsTrunc} and
Dimitrov--Xu~\cite{DX} are essential prior
work on this problem.

The algebraic signs, all-order leading
coefficients and first-order cross identity
are independently tested in
\path{certificates/xi_mode_pair_weighted_exact.py}.
The Lean module
\path{XiKernel/WeightedModePairs.lean}
checks the finite polynomial and
cancellation gates; it does not
formalize the infinite theta modular
identity, Fourier integration by parts,
or the analytic asymptotic error term.
